\section{Frobenius comparison}
\label{sec:frobenius}

The scalar estimate of Lemma~\ref{lemma:scalar-bound} and the
projective-bundle estimate of Proposition~\ref{prop:large-seshadri}
will now give the contradiction needed for
Theorem~\ref{thm:terminal-big}.  Their incompatibility is expressed
by the following theorem about fixed complex data.  It is the
finite-data Frobenius theorem of \cite[Theorem~9.1]{LA}.
We give its proof, including the bounds that
must be uniform when the residual characteristic varies.

The argument compares two orders of vanishing of one determinant.
Ordinary jets make the determinant nonzero after reduction to positive
characteristic. A small polarization bounds the rank on the diagonal,
forcing high vanishing at each diagonal point. A fixed movable curve on a point
blowup bounds the vanishing of every section of each determinant
factor and gives the opposite inequality.

\subsection{The finite-data statement}

For a vector bundle on a smooth projective curve, its slope is
degree divided by rank. The least and greatest Harder--Narasimhan
slopes are denoted by $\mu_{\min}$ and $\mu_{\max}$.
For a line bundle on a smooth variety, to generate jets through
order $a$ at a point means surjectivity to its stalk modulo
the $(a+1)$-st power of the maximal ideal.

\begin{theorem}[Finite-data Frobenius incompatibility]
\label{thm:finite-data}
Let $W$ be a smooth connected projective complex variety of
dimension $n\ge2$. Let $D,H$ be rational Cartier divisors, with
$H$ ample, let $S$ be an integral Cartier divisor, and fix an
integer $q>0$ and real numbers $r>1$, $\epsilon>0$. Suppose
\begin{equation}\label{fr:intersection-bounds}
 \begin{gathered}
 H^n\le(2\epsilon)^n,\qquad
 K_WH^{n-1}\le2H^n/r,\qquad
 (2D+qS)H^{n-1}\le4H^n,\\
 (14\epsilon)^n<\tfrac14,\qquad 80\epsilon<\tfrac34.
 \end{gathered}
\end{equation}
Assume that the following fixed data exist.
\begin{enumerate}
\renewcommand{\labelenumi}{\textup{(\roman{enumi})}}
\item A smooth integral complete-intersection flag from $W$ to
a curve $C$, whose successive divisor classes are $h_i=l_iH$
for fixed positive integers $l_i$ for which $l_iH$ is Cartier, with
$\mu_{\min}(\Omega_W^1|_C)\ge0$.
\item On the projective bundle
\[
 Z=\PP\bigl(\mathcal O_W(S)_1\oplus
                         \mathcal O_W(S)_2\bigr)\longrightarrow W\times W,
 \qquad \xi=c_1(\mathcal O_Z(1)),\qquad M=D_1+D_2+q\xi,
\]
use the convention that the direct image of $\mathcal O_Z(a\xi)$
is the $a$-th symmetric power of the displayed sum for $a\ge0$;
subscripts indicate pullback from the corresponding factor of
$W\times W$. The two axes are the sections defined by the two
summands. A smooth point $z_*$ in their complement and finitely many
sections of $k_0M$, for some integer $k_0>0$ with $k_0D$ Cartier,
generate jets through order $(2n+2)k_0$ at $z_*$.
\item A point $x\in W$, its blowup
$\pi_x:\widehat W\to W$ with exceptional divisor $J$, and a
curve class
\[
 \gamma=f_*(A_1\cdots A_{n-1}),
\]
where $f:V\to\widehat W$ is one fixed birational morphism,
$V$ is smooth integral projective, and the $A_i$ are ample
integral Cartier divisors, such that $J\cdot\gamma>0$ and
\begin{equation}\label{fr:movable-input}
 \pi_x^*(D+K_W/2+\lambda S)\cdot\gamma
 \le40\epsilon J\cdot\gamma\qquad(0\le\lambda\le q).
\end{equation}
It suffices to check the two endpoints in this affine inequality.
\end{enumerate}
These data cannot coexist.
\end{theorem}

The point $x$ is independent of the two base coordinates of $z_*$.
No nefness is required of $D,S$, or $K_W$, and the theorem does not
assume that $K_W$ is pseudo-effective. Its hypotheses specify the
flag and the movable curve class themselves. All these data will
be fixed before the residual characteristic tends to infinity.

\subsection{Reduction and Frobenius jets}
\label{fr:reduction}

We begin the proof of Theorem~\ref{thm:finite-data} by spreading
its hypotheses and converting the fixed ordinary jets
into jets of order proportional to the residual characteristic.
Put
$\Lambda=\prod_{i=1}^{n-1}l_i$. The flag satisfies
\begin{equation}\label{fr:flag-degrees}
 [C]=\Lambda H^{n-1},\qquad h_i\cdot C=l_i\Lambda H^n.
\end{equation}
Choose one sufficiently ample integral Cartier divisor $T_0$ so
that every $T_0+aS$, for integers $0\le a\le(k_0-1)q$, has a section
nonvanishing at each of the two base coordinates of $z_*$.
Fix these finitely many sections. All models, morphisms, divisors,
points, the flag, the complete-intersection witness for $\gamma$,
and the jet and filler sections spread over an integral finitely
generated $\Z$-algebra of characteristic zero. Shrink its
spectrum so that the fibers and flag are smooth projective and
geometrically integral, the birational morphisms remain birational,
the fixed ample bundles remain ample, and the finite jet
surjection and nonvanishings persist. Birationality is preserved
by spreading an isomorphism between dense opens. All displayed
intersection numbers are constant.

We also preserve $\mu_{\min}(\Omega_W^1|_C)\ge0$.
One way is to spread its characteristic-zero Harder--Narasimhan
filtration, make all its factors locally free, and use openness
of their semistability on curves; their degrees and ranks are
fixed. Equivalently, a single relative ample twist globally
generates the curve bundle. Every quotient of rank $a$ then has
degree bounded below by $-ad$ for a fixed integer $d$.
A negative-degree quotient therefore has one of finitely many ranks
and degrees. The proper relative Quot schemes for those Hilbert
polynomials have images missing the generic point: after removal
of torsion, a negative-degree coherent quotient there would
give a negative-degree vector-bundle quotient. Removing those
finitely many images proves the same openness assertion directly.

This base has closed points in arbitrarily large characteristics.
Take algebraic closures of their residue fields and exclude
denominator primes. Retain the notation for the reductions.
The curve $\gamma$ continues to pair nonnegatively with every
effective divisor: pull that divisor back under the fixed
birational morphism and intersect with its fixed ample divisors.
This tests effective divisors on the reduction itself. No
specialization of a characteristic-zero effective cone is used.

\begin{lemma}\label{fr:positive-characteristic-jets}
For a sufficiently large residual characteristic $p$, set
\[
 N_p=\lfloor p/k_0\rfloor,\qquad B_p=k_0N_pD+T_0.
\]
The sections of $pq\xi+B_{p,1}+B_{p,2}$ generate jets through order
$(2n+2)k_0N_p$ at $z_*$.  In particular, they surject onto the quotient
of its local ring by the $p$-th powers of all regular parameters.
\end{lemma}

\begin{proof}
Multiply $N_p$ copies of the fixed jet system.  Every monomial of
total degree at most $(2n+2)k_0N_p$ is a product of $N_p$ monomials
of degree at most $(2n+2)k_0$.  Taking sections with those leading
monomials gives a triangular system, ordered by total degree.
Starting at degree zero and removing the error in each successive
degree proves the required jet surjection. This argument takes
place in the local ring and never divides by factorials.

Put $d_p=(p-k_0N_p)q$, so $0\le d_p\le(k_0-1)q$.
Among the sections fixed before reduction, choose
\[
 u_0\in H^0(W,\OO_W(T_0)),\qquad
 v_{d_p}\in H^0(W,\OO_W(T_0+d_pS))
\]
nonzero at the first and second base coordinates of $z_*$,
respectively. Their tensor product belongs to the summand of
first-slot weight zero in the projective-bundle formula; its
fiber monomial is the $d_p$-th power of the second coordinate.
It defines a section of
$d_p\xi+T_{0,1}+T_{0,2}$ nonzero at $z_*$, since $z_*$ lies
off both axes. Multiplying the product jet system by this section
changes its divisor from $N_pk_0M$ to
$pq\xi+B_{p,1}+B_{p,2}$. Multiplication is invertible on the
local jet algebra. Only the finitely many previously fixed
values of $d_p$ occur.

Since $\dim Z=2n+1$, the
quotient by the $p$-th powers of regular parameters has top total
degree $(2n+1)(p-1)$.  For large $p$ this is at most
$(2n+2)k_0N_p$. The underlying ordinary/Frobenius ideal comparison is
also recorded in \cite[proof of Proposition~2.12, Equation~(2.8)]{MustataSchwede}.
\end{proof}

\subsection{Full rank from the two-slot jets}
\label{fr:full-rank}

Let $F:W\to W'$ be relative Frobenius to the base-field twist.
It is finite flat because $W$ is smooth.  Write $S'$ for the twist
of $S$, so $F^*S'=pS$, and define
\[
 \mathcal E=F_*\OO_W(B_p),\qquad s=\rk\mathcal E=p^n,
 \qquad R=s^2.
\]
For integers $0\le a\le q$, put
$U_a=H^0(W,\OO_W(B_p+paS))$.
Projection formula gives an evaluation map on $W'\times W'$:
\begin{equation}\label{fr:two-slot-evaluation}
 \bigoplus_{\substack{a+b=q\\a,b\ge0}}
 U_a\otimes U_b\otimes
 \OO(-aS'_1-bS'_2)
 \longrightarrow \mathcal E\boxtimes\mathcal E.
\end{equation}

\begin{lemma}\label{fr:generic-rank}
The map in \eqref{fr:two-slot-evaluation} has generic rank $R$.
\end{lemma}

\begin{proof}
Trivialize the two summands defining $Z$ near the selected base
points.  Let $z$ be the torus coordinate, whose value at $z_*$ is
$z_0\ne0$.  By the projective-bundle formula, the sections in
Lemma~\ref{fr:positive-characteristic-jets} are weight polynomials in $z$
with weight-$j$ coefficients in
\[
 H^0(W,B_p+jS)\otimes
 H^0(W,B_p+(pq-j)S),\qquad 0\le j\le pq.
\]
Let $A_*$ be the tensor product of the local algebras of the two
base Frobenius fibers.  The local quotient in
Lemma~\ref{fr:positive-characteristic-jets} is
\[
 A_*[z]/(z^p-z_0^p),
\]
free over $A_*$ with basis $1,z,\ldots,z^{p-1}$.  Project to the
coefficient of $1$ in this basis.  Exactly the weights $j=pa$
survive, with multipliers $z_0^{pa}$.  Their coefficients are the
values of the corresponding summands of
\eqref{fr:two-slot-evaluation}, in the chosen frames and the induced
twisted frames.  As the full jet map is surjective, these coefficients
span $A_*$, of dimension $p^{2n}=R$.  Thus the evaluation has full
rank at this pair of base points, and hence generically.  The
calculation uses a basis over the possibly nonreduced algebra $A_*$,
so reducedness of the Frobenius fibers is unnecessary.
\end{proof}

The jet calculation at $z_*$ proves generic full rank; the diagonal
estimate below will be tested at $(x',x')$, obtained from the
separately fixed point $x$.

\subsection{Sections on the Frobenius diagonal}
\label{fr:filtration}

The Frobenius fiber product has the cartesian square
\begin{equation}\label{fr:frobenius-square}
\begin{tikzcd}
 \Delta_F=W\times_{W'}W \arrow[r,"\mathrm{pr}_2"] \arrow[d,"\mathrm{pr}_1"'] & W \arrow[d,"F"]\\
 W \arrow[r,"F"'] & W'.
\end{tikzcd}
\end{equation}
Its reduced subscheme is the diagonal $W$.  Put
\[
 \mathcal L_p=
 \OO_{\Delta_F}(B_{p,1}+B_{p,2}+pqS_1).
\]
Write $h=F\circ\mathrm{pr}_1=F\circ\mathrm{pr}_2:
\Delta_F\to W'$. On this fiber product,
\[
 \OO(pS_1)\simeq h^*\OO_{W'}(S')\simeq\OO(pS_2).
\]
Consequently every pair of weights $a+b=q$ gives the same line bundle:
\[
 \OO_{\Delta_F}(B_{p,1}+B_{p,2}+paS_1+pbS_2)
 \simeq\mathcal L_p.
\]
On the diagonal of $W'\times W'$, every source twist of
\eqref{fr:two-slot-evaluation} becomes $-qS'$.
Finite base change and projection formula identify the twisted target as
\[
 (\mathcal E\otimes\mathcal E)(qS')\simeq h_*\mathcal L_p.
\]
The product sections defining the columns restrict to global sections
of this one bundle $\mathcal L_p$. Their values at a diagonal point
therefore lie in the image of $H^0(\Delta_F,\mathcal L_p)$ in the
corresponding fiber of $h_*\mathcal L_p$. This proves that the rank
at every diagonal point is at most
\begin{equation}\label{fr:diagonal-rank-upper}
 h^0(\Delta_F,\mathcal L_p).
\end{equation}

For a rank-$n$ bundle $V$ in characteristic $p$, denote by $A_j(V)$
the degree-$j$ part of its symmetric algebra modulo the $p$-th powers
of local generators.  This construction is independent of frame:
the $p$-th power of a linear combination is the sum of the $p$-th
powers in characteristic $p$.  Each $A_j(V)$ is locally free and a
quotient of $V^{\otimes j}$.  Set
\[
 u=n(p-1),\qquad e_j=\rk A_j(V),\qquad
 \sum_{j=0}^u e_j=p^n.
\]
Filtering $\mathcal L_p$ by powers of the ideal of the reduced
diagonal in $\Delta_F$ gives the graded bundles
\begin{equation}\label{fr:diagonal-graded}
 \mathcal G_j=\OO_W(2B_p+pqS)\otimes A_j(\Omega_W^1),
 \qquad 0\le j\le u.
\end{equation}
Indeed, in smooth local coordinates the algebra is generated by
parameter differences with their $p$-th powers zero, and its
degree-one conormal is $\Omega_W^1$.  This coordinate calculation
can be made étale locally or in completions and identifies the
global associated graded.  This is the diagonal-ideal form of the canonical Frobenius filtration
\cite[Corollary~3.5]{KS}. The same calculation, using
only the bundle from the second factor, filters $F^*\mathcal E$ with pieces
\cite[Theorem~3.7]{Sun}
\begin{equation}\label{fr:Frob-pullback-graded}
 \OO_W(B_p)\otimes A_j(\Omega_W^1).
\end{equation}

Complementary multiplication is a perfect pairing, giving
\begin{equation}\label{fr:socle-pairing}
 A_{u-j}(V)\simeq A_j(V)^\vee\otimes(\det V)^{p-1}.
\end{equation}
On monomials, each exponent vector pairs with its complement to
$(p-1,\ldots,p-1)$.  The top line transforms by the character
$(\det V)^{p-1}$: this is immediate on diagonal matrices and hence
 identifies the character of $\mathrm{GL}_n$ on that line.

\subsection{A uniform bound for the diagonal rank}
\label{fr:rank-bound}

We bound the sections of every graded bundle in
\eqref{fr:diagonal-graded} by its rank times the same scalar
polynomial in $p$. This uniformity lets us sum the ranks of the
graded pieces without an additional factor for their number.
We begin with a slope bound on the fixed curve; here slope means
ordinary degree divided by rank.

\begin{lemma}\label{fr:truncated-slope}
There is a constant $c\ge0$, independent of the sufficiently large
residual characteristic $p$ and of $j$, such that for
$V=\Omega_W^1|_C$ one has
\[
 \mu_{\max}(A_j(V))\le(p-1)K_W\cdot C+nc.
\]
\end{lemma}

\begin{proof}
Write $F_C$ for absolute Frobenius of $C$.  Langer's instability
estimate \cite[Corollary~2.5]{Langer} gives
\begin{equation}\label{fr:Langer-min}
 L_{\min}(V):=
 \lim_{e\to\infty}p^{-e}\mu_{\min}(F_C^{e*}V)
 \ge-\frac{c}{p-1}.
\end{equation}
To make the uniformity explicit, let $g$ be the fixed genus of $C$.
One can choose a nef line bundle $A_C$ of degree at most
$d_g=\max\{4g-2,0\}$ such that $T_C\otimes A_C$ is globally
generated; for $g\le1$, take $A_C=\OO_C$. Langer's bound is
\[
 \mu_{\min}(V)-L_{\min}(V)
 \le\frac{(\rk V-1)\deg A_C}{p-1}.
\]
Since $\rk V=n$ and $\mu_{\min}(V)\ge0$, we may take
$c=(n-1)d_g$, independently of the reduction.

We also need, for every integer $i\ge0$,
\begin{equation}\label{fr:tensor-min}
 \mu_{\min}(V^{\otimes i})\ge iL_{\min}(V).
\end{equation}
Fix $p$ and $i$.  Choose a line bundle $A_e$ of degree
$-\mu_{\min}(F_C^{e*}V)+O(1)$, rounded up with a fixed
genus-dependent additive constant, so that
$F_C^{e*}V\otimes A_e$ is globally generated.  This follows from
Serre duality: after subtracting any point, its minimum slope can
be made greater than $2g(C)-2$, which annihilates $H^1$ and makes
evaluation at that point surjective.  For every vector-bundle
quotient $Q$ of $V^{\otimes i}$, the bundle
$F_C^{e*}Q\otimes A_e^{\otimes i}$ is then globally generated.
Consequently
\[
 p^e\mu(Q)+i\deg A_e\ge0.
\]
Divide by $p^e$ and let $e\to\infty$ to prove
\eqref{fr:tensor-min}.  This limit is taken for each fixed $p,i$;
there is no interchange of Frobenius-iteration and characteristic
limits.

Since $A_{u-j}(V)$ is a quotient of $V^{\otimes(u-j)}$, the perfect
pairing \eqref{fr:socle-pairing} yields
\begin{align*}
 \mu_{\max}(A_j(V))
 &= (p-1)\deg\det V-\mu_{\min}(A_{u-j}(V))\\
 &\le(p-1)K_W\cdot C+\frac{(u-j)c}{p-1}\\
 &\le(p-1)K_W\cdot C+nc.
\end{align*}
\end{proof}

The divisors $B_p-pD=T_0-(p-k_0N_p)D$ range over a fixed finite
list. Combining~\eqref{fr:intersection-bounds},
\eqref{fr:flag-degrees}, and Lemma~\ref{fr:truncated-slope}
therefore gives, uniformly for $0\le j\le u$,
\begin{align}
 \mu_{\max}(\mathcal G_j|_C)
 &\le p(4+2/r)\Lambda H^n+O(1)\notag\\
 &\le M_p:=7p\Lambda H^n+c_0,\label{fr:common-slope-bound}
\end{align}
where $c_0\ge0$ is fixed.  The coefficient $7$ provides harmless
room in this common bound.

\begin{lemma}\label{fr:diagonal-sections}
For all sufficiently large $p$,
\[
 h^0(\Delta_F,\mathcal L_p)
 \le R\bigl(7^nH^n+O(1/p)\bigr)<R/4.
\]
The error constant is independent of the filtration index $j$.
\end{lemma}

\begin{proof}
Write the fixed flag as
\[
 W=W_n\supset W_{n-1}\supset\cdots\supset W_1=C,
 \qquad W_{n-k}\in|h_k|_{W_{n-k+1}}
 \quad(1\le k\le n-1).
\]
At step $k$, fix the preceding nonnegative integers
$b_1,\ldots,b_{k-1}$ and put
\[
 \mathcal F_k=
 \mathcal G_j|_{W_{n-k+1}}
 \Bigl(-\sum_{i<k}b_i h_i\Bigr).
\]
For each $b_k\ge0$, the divisor restriction sequence is
\[
 0\longrightarrow\mathcal F_k(-(b_k+1)h_k)
 \longrightarrow\mathcal F_k(-b_kh_k)
 \longrightarrow
 \mathcal F_k(-b_kh_k)|_{W_{n-k}}
 \longrightarrow0.
\]
Iterating it for $b_k=0,1,\ldots$ and then proceeding to the next
member of the flag bounds the section space by
\begin{equation}\label{fr:flag-lattice-sum}
 h^0(W,\mathcal G_j)\le
 \sum_{b_1,\ldots,b_{n-1}\ge0}
 h^0\left(C,\mathcal G_j|_C\Bigl(-\sum_i b_i h_i|_C\Bigr)\right).
\end{equation}
At each stage the remainder is zero after sufficiently negative
ample twisting.  Its stopping index need not be uniform in $j,p$
or in the preceding twists: enlarging the nonnegative finite sums
to the displayed lattice sum preserves the inequality.

A summand in \eqref{fr:flag-lattice-sum} vanishes when
$\sum_i b_i(h_i\cdot C)>M_p$, because then its maximum slope is
negative.  Each nonzero summand has at most $e_j(M_p+1)$ sections.
Indeed, evaluation at $\lfloor M_p\rfloor+1$ distinct points is injective;
the kernel has negative maximum slope after subtracting those
points.  By \eqref{fr:flag-degrees} it follows that
\begin{align}
 h^0(W,\mathcal G_j)
 &\le e_j(M_p+1)
       \prod_{i=1}^{n-1}
       \left(1+\frac{M_p}{l_i\Lambda H^n}\right)\notag\\
 &=e_jp^n\bigl(7^nH^n+O(1/p)\bigr).
 \label{fr:common-flag-polynomial}
\end{align}
The leading coefficient follows from
$\Lambda=\prod_i l_i$.  More importantly, the whole scalar
polynomial on the first line is the same for every $j$.

Sum over the filtration in \eqref{fr:diagonal-graded}.  The ranks
satisfy $\sum_j e_j=p^n$, so \eqref{fr:common-flag-polynomial} gives
the asserted bound with $R=p^{2n}$.  There is no extra factor for
the number of graded pieces.  Finally
\[
 7^nH^n\le(14\epsilon)^n<1/4
\]
by \eqref{fr:intersection-bounds}, proving the strict inequality
for all large $p$.
\end{proof}

\subsection{The two orders of the determinant}
\label{fr:determinant}

Choose $R$ generically independent columns of
\eqref{fr:two-slot-evaluation} from bases of its source vector spaces.
Their determinant is a nonzero section
\[
 0\ne\sigma\in H^0(W'\times W',\mathcal Q_1\boxtimes\mathcal Q_2)
\]
of an actual external-product line bundle. To specify its factors,
let $m_i$ be the sum of the weights of the $R$ chosen columns in
slot $i$. The target determinant is
$(\det\mathcal E)^s\boxtimes(\det\mathcal E)^s$, and the
chosen source determinant is
$\OO(-m_1S')\boxtimes\OO(-m_2S')$. Thus
\[
 \mathcal Q_i=(\det\mathcal E)^{\otimes s}
                 \otimes\OO_{W'}(m_iS'),\qquad
 \lambda_i=m_i/R.
\]
Every column weight lies between $0$ and $q$, so
$0\le\lambda_i\le q$ independently of the columns chosen.

Under the numerical identification with the base-field twist,
\begin{equation}\label{fr:slot-determinant-class}
 \frac{c_1(\mathcal Q_i)}R
 =\frac{c_1(\mathcal E)}s+\lambda_iS
 =\frac{B_p}p+\frac{p-1}{2p}K_W+\lambda_iS,
 \qquad 0\le\lambda_i\le q.
\end{equation}
The Frobenius first-Chern-class identity is also
\cite[Lemma~4.2]{Sun}; it is an identity of rational divisor classes.  For the second equality, use
\eqref{fr:Frob-pullback-graded}.  Complementary degrees in
\eqref{fr:socle-pairing} have total first Chern class
$e_j(p-1)K_W$.  Summing gives
\[
 \sum_jc_1(A_j(\Omega_W^1))=\frac{s(p-1)}2K_W.
\]
Thus $c_1(F^*\mathcal E)=sB_p+s(p-1)K_W/2$, and Frobenius pullback of
twisted divisor classes multiplies by $p$, as required.

For $0\le\lambda\le q$, the fixed movable inequality gives
\[
 (D+K_W/2+\lambda S)\cdot\gamma\le40\epsilon J\cdot\gamma,
\]
where pullback by $\pi_x$ is understood. This concerns only
intersections of fixed divisors and the fixed complete-intersection
witness, so it remains true on the reduction. The class in
\eqref{fr:slot-determinant-class} differs from its
left-hand divisor at $\lambda=\lambda_i$ by
\[
 \frac1p\bigl(B_p-pD-K_W/2\bigr).
\]
Its numerator belongs to the fixed finite list
$T_0-aD-K_W/2$, $a\in\{0,\ldots,k_0-1\}$. The resulting error after
intersection with $\gamma$ is therefore $O(p^{-1})$, uniformly
over the selected determinant columns.

Let $x'$ be the twisted point and let $0\ne\tau_i\in
H^0(W',\mathcal Q_i)$.  The strict transform of its effective divisor on
the blowup at $x'$ pairs nonnegatively with the twist of $\gamma$.
Since $J\cdot\gamma>0$, it follows that
\begin{equation}\label{fr:slot-order}
 \frac{\ord_{x'}(\tau_i)}R\le40\epsilon+O(1/p).
\end{equation}
The constant is uniform in the columns and in $\tau_i$.  This
argument tests sections on the reduction itself; it does not
require them to lift to characteristic zero.

K\"unneth identifies the space containing $\sigma$ with
$H^0(W',\mathcal Q_1)\otimes H^0(W',\mathcal Q_2)$.  Choose bases adapted to the
orders of vanishing at $x'$ in each factor.  In each degree their
leading terms are linearly independent.  Tensor products of these
leading terms remain independent in each bidegree of the two
disjoint sets of parameters; terms of different bidegrees cannot
cancel.  Every nonzero tensor consequently has order at most the
sum of the two maximum basis orders.  Applying
\eqref{fr:slot-order}, we obtain
\begin{equation}\label{fr:determinant-order-upper}
 \ord_{(x',x')}\sigma\le R\bigl(80\epsilon+O(1/p)\bigr).
\end{equation}

On the other hand, \eqref{fr:diagonal-rank-upper} and
Lemma~\ref{fr:diagonal-sections} bound the matrix rank at $(x',x')$
by a number strictly smaller
than $R/4$.  Trivialize its line bundles near that point.  Invertible
constant row operations make more than $3R/4$ rows vanish at the
point.  Each entry of those rows belongs to the maximal ideal, so
every term of the determinant has order at least $3R/4$.  Therefore
\begin{equation}\label{fr:determinant-order-lower}
 \ord_{(x',x')}\sigma\ge3R/4.
\end{equation}
Equations \eqref{fr:determinant-order-upper} and
\eqref{fr:determinant-order-lower} contradict
$80\epsilon<3/4$ for sufficiently large $p$.


This contradiction proves Theorem~\ref{thm:finite-data}.
The full-rank calculation used $z_*$, while the diagonal estimate
holds at every point and was tested at the separately chosen $x$.
The errors were controlled by one fixed finite list of divisor
classes and one fixed flag polynomial. Thus their constants do
not depend on $p$, the filtration index, or the chosen columns.

\subsection{Application to the terminal case}
\label{fr:application}

We return to Theorem~\ref{thm:terminal-big}. Thus $X$ is a projective
terminal $\Q$-factorial variety of dimension $n\ge2$ with
$K_X\sim0$, and $L$ is a nef Cartier divisor of full nef dimension.
Assume that $L$ is not big. Recall the ample rational divisors
$P=r(L+tA)$, where $A$ is a fixed ample Cartier divisor and $t>0$
tends to zero, with
\[
 r\longrightarrow\infty,\qquad P^n\longrightarrow\epsilon^n,
 \qquad \rho=(P^n)^{1/n}.
\]
The positive rational number $\epsilon$ was chosen so that
$(14\epsilon)^n<1/4$ and $80\epsilon<3/4$. Choose the integer $q$
supplied by Proposition~\ref{prop:large-seshadri}, and then choose
$t$ sufficiently small that this proposition and
Lemma~\ref{lemma:scalar-bound} apply, that $r>q+1$, and that
$P^n<(2\epsilon)^n$. We now fix $t$, and therefore also $P,r,\rho$
and the projective bundle in~\eqref{eq:bundle}. Every remaining
choice is made over $\C$ before reducing to positive characteristic.

\paragraph{The divisors and the flag.}
Choose a smooth projective resolution $w:W\to X$ and set
\[
 D=w^*P,\qquad S=w^*L,\qquad H_0=w^*P.
\]
The divisor $S$ is integral Cartier, and $H_0$ is nef and big.
Since $X$ is terminal and $K_X\sim0$, there is an integral
effective $w$-exceptional divisor $E$ with $K_W\sim E$. The
projection formula and the nefness of $P-rL$ give
\[
 K_WH_0^{n-1}=0,\qquad
 (2D+qS)H_0^{n-1}\le(2+q/r)H_0^n.
\]
At $H_0$, the inequalities
$K_WH_0^{n-1}<2H_0^n/r$ and
$(2D+qS)H_0^{n-1}<4H_0^n$ are strict. Thus, by continuity of
intersection products, a sufficiently small rational ample
perturbation $H=H_0+\eta A_W$, with $A_W$ an ample Cartier divisor
on $W$ and $\eta>0$, satisfies
\[
 H^n\le(2\epsilon)^n,\qquad
 K_WH^{n-1}\le2H^n/r,\qquad
 (2D+qS)H^{n-1}\le4H^n.
\]
This gives~\eqref{fr:intersection-bounds}.

The canonical divisor of $W$ is pseudo-effective. Generic
semipositivity of the cotangent bundle, originating in \cite{Miyaoka}
and used here in the form \cite[Theorem~2.1]{CP}, together with
restriction to sufficiently general complete-intersection curves
\cite[Theorem~6.1]{MR}, gives a smooth integral flag with
successive divisor classes $l_iH$, for sufficiently large divisible
integers $l_i$, such that
\[
 \mu_{\min}(\Omega_W^1|_C)\ge0.
\]
This is the flag in Hypothesis~\textup{(i)} of
Theorem~\ref{thm:finite-data}.

\paragraph{The finite jet system.}
On the bundle of~\eqref{eq:bundle}, write
$R=P_1+P_2+q\xi$. The divisor $R$ is ample: $\xi$ is nef and
relatively ample, and $P_1+P_2$ is the pullback of an ample
divisor on $X\times X$. Proposition~\ref{prop:large-seshadri} supplies
a smooth point $z_*$ in the complement of the two axes with
\[
 \varepsilon(R;z_*)>2n+2.
\]
The point can be chosen over the open set where $w$ is an
isomorphism. Choose a rational number $a$ strictly between
$2n+2$ and $\varepsilon(R;z_*)$. On the blowup
$b:\widetilde Z\to Z$ at $z_*$, with exceptional divisor $G$,
the rational divisor $b^*R-aG$ is ample. For sufficiently large
divisible $k_0$, Serre vanishing gives
\[
 H^1\bigl(\widetilde Z,
          \OO_{\widetilde Z}(k_0b^*R-k_0aG)\bigr)=0.
\]
Since $z_*$ is smooth, the sections on the thickened exceptional
divisor $k_0aG$ identify with the jet quotient at $z_*$ modulo
$\mathfrak m_{z_*}^{k_0a}$, in a local frame for $k_0R$.
The restriction sequence therefore shows that sections of
$k_0R$ generate jets through order $k_0a-1$, and hence through
order $(2n+2)k_0$. Increase the divisibility of $k_0$ to make
$k_0P$ Cartier. Pulling these sections to
\[
 Z_W=\PP\bigl(\OO_W(S)_1\oplus\OO_W(S)_2\bigr)
       \longrightarrow W\times W
\]
preserves the jets at the corresponding point, because the induced
map $Z_W\to Z$ is an isomorphism near that point. The pulled-back divisor is
$D_1+D_2+q\xi$. We have obtained the finite system required in
Hypothesis~\textup{(ii)}.

\paragraph{The movable curve class.}
Choose a very general point $x\in W$ outside the exceptional locus,
with image in the locus where Lemma~\ref{lemma:scalar-bound}
holds. Let $\pi_x:\widehat W\to W$ be its blowup and $J$ its
exceptional divisor. The rational divisor
\[
 D^+=2D+E
\]
is big. For every sufficiently divisible positive integer $m$,
normality of $X$ and exceptionality of $E$ give
$w_*\OO_W(mE)=\OO_X$. Thus every section of $mD^+$ comes from a
section of $2mP$, multiplied by the section of the fixed divisor
$mE$. Since $x\notin\Supp E$, Lemma~\ref{lemma:scalar-bound}
implies
\begin{equation}\label{fr:application-order}
 \ord_x(s)/m\le8\rho<16\epsilon
 \qquad\bigl(0\ne s\in H^0(W,\OO_W(mD^+))\bigr).
\end{equation}

It follows that
\begin{equation}\label{fr:application-not-pe}
 \pi_x^*D^+-40\epsilon J\quad\text{is not pseudo-effective}.
\end{equation}
Indeed, if it were pseudo-effective, a convex combination with
the big divisor $\pi_x^*D^+$ would make
$\pi_x^*D^+-cJ$ big for any rational $c$ with
$16\epsilon<c<40\epsilon$. A nonzero section of a sufficiently
divisible multiple would then give a section violating
\eqref{fr:application-order}.

By the duality between pseudo-effective divisors and strongly
movable curves~\cite{BDPP}, the strict separation
in~\eqref{fr:application-not-pe} is detected by a class
\[
 \gamma=f_*(A_1\cdots A_{n-1}),
 \qquad f:V\longrightarrow\widehat W,
\]
where $V$ is smooth projective, $f$ is birational, and the $A_i$
are ample integral Cartier divisors. To obtain this form, first
choose a generating strongly movable class with negative pairing,
then approximate its ample real classes by rational ones and
clear denominators. We have
\[
 (\pi_x^*D^+)\cdot\gamma<40\epsilon J\cdot\gamma.
\]
The left side is nonnegative, because $\pi_x^*D^+$ is
pseudo-effective, so $J\cdot\gamma>0$.

For every $0\le\lambda\le q$,
\[
 D^+-(D+K_W/2+\lambda S)
 \sim_{\R}w^*(P-\lambda L)+E/2
\]
is pseudo-effective: $P-\lambda L=(P-rL)+(r-\lambda)L$ is nef
and $E$ is effective. Pulling back to $\widehat W$ and pairing
with $\gamma$ therefore gives
\[
 \pi_x^*(D+K_W/2+\lambda S)\cdot\gamma
 \le(\pi_x^*D^+)\cdot\gamma
 <40\epsilon J\cdot\gamma.
\]
This proves Hypothesis~\textup{(iii)}. All the hypotheses of
Theorem~\ref{thm:finite-data} now hold for one fixed collection of
complex data, a contradiction. Hence $L$ is big, completing the
proof of Theorem~\ref{thm:terminal-big}.
